The previous section committed us to grounding on the item side. Before
presenting the design, we fix what it must deliver. The design must
\begin{itemize}
\item relate attributes to prototypes in an intrinsic and interpretable
way, readable from the trained parameters rather than recovered after
training,
\item keep the ProtoMF paradigm, so that whatever changes can be
attributed to the one modification,
\item make the attributes available as a source of predictive power,
not only as labels for interpretation.
\end{itemize}

We introduce our approach as feature-composed I-ProtoMF, in short
fI-ProtoMF. It differs from I-ProtoMF~\cite{melchiorre2022protomf} in
exactly one place. The host draws the embedding $\mathbf{t}$ of every
item from a lookup table with one free learned row per item. fI-ProtoMF
replaces this table by a composition in the manner of
LightFM~\cite{kula2015lightfm},
\begin{equation}
\mathbf{t} \,=\, \sum_{f \in \mathcal{F}_t} \mathbf{e}_f
\;\,\bigl(+\; \mathbf{e}_{\mathrm{ID}(t)}\bigr),
\label{eq:fi-composition}
\end{equation}
where $\mathcal{F}_t$ holds the feature values recorded for the item,
each $\mathbf{e}_f \in \mathbb{R}^d$ is a learned vector, and
$\mathbf{e}_{\mathrm{ID}(t)}$ is an optional per-item row discussed
below. Everything downstream remains the host. The $L^t$ item prototypes
$\mathbf{p}^t_l$ stay free learned vectors, the item representation
$\mathbf{t}^* \in \mathbb{R}^{L^t}$ is still the vector of similarities
$\mathrm{sim}(\mathbf{t}, \mathbf{p}^t_l)$, the user is still a free
vector $\mathbf{u} \in \mathbb{R}^{L^t}$, and the score
$\text{I-score}(u,t) = \mathbf{u}^\top \mathbf{t}^*$ is trained with the
host's loss and regularizers (Section~\ref{sec:bg-protomf}). The model
name keeps the f of feature rather than attribute because the
composition sums whatever encoded rows it is given, the ID row being the
first example, and choosing attribute-derived rows is what makes the
resulting model interpretable.

The composition turns attribute values into learned words. Every
distinct value that appears in the catalogue, the colour group black or
the garment group jersey, owns exactly one row $\mathbf{e}_f$, shared by
all items that carry it. An item is then spelled from these words. A
plain black t-shirt sums the rows of its recorded values for department,
product type, section, colour group, and pattern, and a second shirt
that differs only in colour shares four of its five summands. The host's
item table, one independent row per item, becomes a small vocabulary of
attribute rows from which all item embeddings are assembled. The rows
are trained by the collaborative signal like any other parameter, but
each row receives the gradients of every interaction with every item
that carries its value, so what a row comes to encode is the
collaborative meaning of one attribute value across the catalogue.

This is the grounding announced in the previous section, stated
precisely. What the composition constrains directly is the item
embedding. Every $\mathbf{t}$ is a sum of rows that carry meaning on
their own. The prototypes are not constrained at all. They remain free
vectors, but they anchor a space that is now populated by attribute
words, and a prototype can therefore be read by comparing it to those
words directly in parameter space, with no detour through training data.
Grounding the prototypes is in this sense shorthand for grounding the
space they anchor, and the read-outs this buys are developed in
Section~\ref{sec:fi-readouts}.

What the composition gains can be stated as one identity. For any
prototype $\mathbf{p}$, the inner product with an item embedding splits
over the summands of Equation~\ref{eq:fi-composition},
\begin{equation}
\langle \mathbf{t}, \mathbf{p} \rangle
\,=\, \sum_{f \in \mathcal{F}_t} \langle \mathbf{e}_f, \mathbf{p} \rangle
\;\,\bigl(+\; \langle \mathbf{e}_{\mathrm{ID}(t)}, \mathbf{p} \rangle\bigr),
\label{eq:fi-decomposition}
\end{equation}
and dividing by the norms turns the item cosine into a norm-weighted
sum of attribute cosines,
\begin{equation}
\cos(\mathbf{t}, \mathbf{p})
\,=\, \sum_{f \in \mathcal{F}_t}
\frac{\lVert \mathbf{e}_f \rVert}{\lVert \mathbf{t} \rVert}\,
\cos(\mathbf{e}_f, \mathbf{p}),
\label{eq:fi-cosine}
\end{equation}
where the ID term, when present, joins the sum like any other row. The
host's shifted similarity adds a constant and changes nothing here. The
host interprets a prototype by ranking all items by
$\cos(\mathbf{t}, \mathbf{p})$ and reading the metadata of the top of
the list, so it shows aggregates of exactly the terms that
Equation~\ref{eq:fi-cosine} names individually. An item is near a
prototype because its sum aligns, and the list alone cannot say whether
one attribute carries the alignment, several share it, or the ID row
does the work. Ranking the rows $\mathbf{e}_f$ by
$\cos(\mathbf{e}_f, \mathbf{p})$ reads the terms themselves, one
meaning at a time, with no detour through the catalogue.

The same equations show what is lost. Every item embedding now lies in
the span of the vocabulary rows, so the item table's capacity drops
from one free direction per item to at most one per attribute value,
and two items that record the same attribute values receive the same
embedding, the same representation, and the same score for every user
unless the ID row separates them. Equation~\ref{eq:fi-cosine} also
carries a caveat for the read-out itself. Its weights
$\lVert \mathbf{e}_f \rVert / \lVert \mathbf{t} \rVert$ are set by the
row norms, so a row with a large norm dominates every item cosine it
enters, and the nearest items of a prototype and its nearest attribute
rows need not agree. A prototype may also align with a direction
between words that no single row captures, in which case the attribute
ranking is flat while the item ranking still concentrates. The two
read-outs are therefore complements rather than substitutes, and
Section~\ref{sec:fi-readouts} uses both.

\paragraph{The optional ID row.} The row $\mathbf{e}_{\mathrm{ID}(t)}$
in Equation~\ref{eq:fi-composition} is a free per-item vector added to
the sum by configuration. It restores per-item capacity, letting an item
deviate from what its attribute values alone would predict, at the price
of an ungrounded summand in every embedding. It also makes the host a
special case. With $\mathcal{F}_t$ empty and the ID row kept, the model
reduces to I-ProtoMF exactly, and our implementation reproduces the host
bit-identically in this configuration. The chapter therefore evaluates
two arms against the host. fI-ProtoMF uses features and the ID row,
and the noid arm uses features alone. What the ID row
contributes, and what its ungrounded share costs, is measured in
Sections~\ref{sec:fi-readouts} and~\ref{sec:fi-results}.
